\section{Experiment 11B: Periodically Calibrated Sparse Descent Certificates}
\label{sec:exp11b}

\subsection{Purpose}

Experiment~11A showed that descent-aware event sizing is extremely effective when the server knows the aggregate gradient, but that exact client-local descent is not a valid substitute under heterogeneous objectives. Experiment~11B therefore asks the next information question:
\begin{equation}
  \boxed{\text{Can occasional dense calibration messages support a much larger stream of safe low-bit coordinate events?}}
\end{equation}
The experiment deliberately retains the normalized full-reset LIF trigger. A threshold crossing still communicates only a coordinate address and sign. What changes is the information maintained by the server for deciding whether the event is globally useful and how large its jump may be.

The study uses the same $N=10$, $d=20$ asynchronous quadratic family as Experiments~10D and 11A. Two geometries are evaluated: the diagonal problem and the strongly rotated isospectral problem from Experiment~10D. The primary campaign uses 40 held-out problems, 2000 wall-clock ticks, a 500-tick tail window, and stochastic-gradient noise $\sigma_g=0.25$.

\subsection{Calibration model and communication accounting}

At calibration time $t_c$, every client evaluates and transmits its full local gradient at the current server model,
\begin{equation}
  g_i^c=\nabla F_i(w^c).
\end{equation}
The server stores the individual vectors and their aggregate center
\begin{equation}
  c^c=\sum_i p_i g_i^c.
\end{equation}
A full calibration costs
\begin{equation}
  B_{\rm cal}=N d\,32=6400
\end{equation}
logical payload bits before packet headers. This recurring calibration traffic, including the initial calibration at $t=0$, is included in every communication result. Static curvature bounds are treated as configured benchmark metadata and are not charged repeatedly.

\subsection{Conservative drift envelope}

For the quadratic local objective,
\begin{equation}
  \nabla F_i(w)-\nabla F_i(w^c)=H_i(w-w^c).
\end{equation}
The primary certificate uses the coordinate-wise deterministic bound
\begin{equation}
  \abs{\nabla_jF_i(w)-g_{i,j}^c}
  \leq
  \sum_k \abs{[H_i]_{jk}}\abs{w_k-w_k^c}.
  \label{eq:11b-componentwise-bound}
\end{equation}
After weighted aggregation, the server obtains
\begin{equation}
  \nabla_jF(w)\in[c_j-r_j,c_j+r_j],
\end{equation}
where
\begin{equation}
  c_j=\sum_i p_i g_{i,j}^c,
  \qquad
  r_j=\sum_i p_i\sum_k\abs{[H_i]_{jk}}\abs{w_k-w_k^c}.
\end{equation}
For an arriving signed coordinate event $s\in\{-1,+1\}$, a guaranteed lower bound on aggregate descent alignment is
\begin{equation}
  \underline a_j=-s c_j-r_j.
\end{equation}
The event is accepted only when $\underline a_j>0$, with jump
\begin{equation}
  q_j=\frac{\underline a_j}{[\bar H]_{jj}}.
  \label{eq:11b-certified-jump}
\end{equation}
Rejected threshold crossings still count as communication and reset their LIF membrane.

\paragraph{Sequential application under coupled geometry.}
A subtle point is essential once $\bar H$ is not diagonal. If several coordinate events arrive in the same packet, applying all individually certified jumps simultaneously introduces Hessian cross-terms that are not covered by independent coordinate descent certificates. Therefore the final Experiment~11B implementation processes certified coordinates \emph{sequentially} at the server. After each accepted coordinate jump, the current model changes before the next candidate in that packet is evaluated. The reported rotated-geometry results below use this sequential-safe implementation. In the diagonal case the distinction disappears because the cross-terms are zero.

A scalar row-norm/Lipschitz envelope was also tested during the smoke stage. It was valid but far too conservative: useful progress required almost continuous recalibration. The componentwise bound in \cref{eq:11b-componentwise-bound} is therefore the primary rule.

\subsection{Controls}

The comparison contains:
\begin{enumerate}
  \item the tuned no-certificate global $q(t)$ event optimizer;
  \item the exact global descent oracle from Experiment~11A;
  \item the periodic conservative certificate in \cref{eq:11b-certified-jump};
  \item a naive stale-calibration rule that ignores $r_j$;
  \item certificate-triggered recalibration with a minimum refresh interval;
  \item an event-time-tightened variant that replaces the firing client's stale contribution by the first-passage gradient estimate plus an empirical uncertainty margin;
  \item an oracle basis-aligned calibrated certificate in the common eigenbasis from Experiment~10D;
  \item EF-TopK and full-precision asynchronous optimization as conventional references.
\end{enumerate}
The first-passage margins are estimated on a separate calibration ensemble and grouped by client compute period. Both 99\% and 99.9\% empirical error quantiles were inspected. These timing-assisted rules are statistical diagnostics, not deterministic certificates.

\subsection{Diagonal geometry: sparse calibration closes the oracle gap}

The diagonal problem is highly favorable because \cref{eq:11b-componentwise-bound} is tight coordinate by coordinate. Every tested periodic certificate, including $C=1000$, reproduces the global-oracle trajectory to numerical precision:
\begin{center}
\begin{tabular}{lrrrr}
\toprule
Method & Tail gap & Whole-run gap & Payload bits & Calibrations \\
\midrule
$q(t)$ baseline & $3.89\times10^{-3}$ & 0.881 & 10378 & 0 \\
Global descent oracle & $\approx0$ & 0.210 & 7812 & 0 \\
Periodic certificate, $C=500$ & $\approx0$ & 0.210 & 39812 & 5 \\
Periodic certificate, $C=1000$ & $\approx0$ & 0.210 & 27012 & 3 \\
\bottomrule
\end{tabular}
\end{center}
Thus only three dense calibrations over 2000 ticks, including the initial one, are sufficient to retain oracle-like descent. The accepted-event fraction is about 1.5\%, and no accepted event is globally harmful.

This is the first practical evidence that Experiment~11A's oracle is reachable from stale but bounded federated information rather than requiring a dense gradient at every event.

\subsection{Rotated geometry: a clear calibration-rate transition}

The strongly rotated isospectral problem is harder because one coordinate displacement changes many gradient components. The same valid uncertainty box therefore expands much faster between calibrations.

The retuned no-certificate reference from Experiment~10D gives
\begin{equation}
  J_{\rm tail}=1.15\times10^{-2},
  \qquad
  J_{\rm whole}=0.294,
  \qquad
  B=1.57\times10^4\;\text{bits}.
\end{equation}
With sequentially processed certified events, the native-coordinate trade-off is
\begin{center}
\begin{tabular}{rrrrr}
\toprule
$C$ & Tail gap & Whole-run gap & Total payload & Calibration payload \\
\midrule
100 & $1.59\times10^{-2}$ & 0.362 & 150163 & 134400 \\
75  & $6.47\times10^{-3}$ & 0.291 & 188359 & 172800 \\
60  & $2.83\times10^{-3}$ & 0.247 & 233012 & 217600 \\
50  & $1.21\times10^{-3}$ & 0.217 & 277748 & 262400 \\
40  & $5.14\times10^{-4}$ & 0.190 & 341674 & 326400 \\
25  & $6.3\times10^{-5}$ & 0.160 & 533597 & 518400 \\
\bottomrule
\end{tabular}
\end{center}
All accepted events in this deterministic-certificate sweep have measured objective change $\Delta F\leq0$; the harmful accepted-event fraction is zero.

The useful transition occurs near $C=75$: it roughly preserves the whole-run cost of the tuned $q(t)$ baseline while cutting the tail error almost in half, but it uses about twelve times more payload. More frequent calibration moves the method steadily toward the global oracle.

\subsection{Comparison with conventional sparse optimization}

On the same strongly rotated benchmark, the selected EF-TopK ($k=4$) reference gives
\begin{equation}
  J_{\rm tail}=7.89\times10^{-3},
  \qquad
  J_{\rm whole}=0.183,
  \qquad
  B=1.0952\times10^6\;\text{payload bits}.
\end{equation}
The native calibrated event method at $C=25$ obtains
\begin{equation}
  J_{\rm tail}=6.3\times10^{-5},
  \qquad
  J_{\rm whole}=0.160,
  \qquad
  B=5.336\times10^5\;\text{bits}.
\end{equation}
Hence this operating point improves whole-run objective by roughly 12\%, reduces the tail gap by more than two orders of magnitude, and uses approximately
\begin{equation}
  \boxed{51\%\text{ fewer logical payload bits}}
\end{equation}
than EF-TopK. With the illustrative packet headers the reduction remains approximately 54\%.

At $C=75$, the method uses only $1.88\times10^5$ bits, about 83\% fewer than EF-TopK, and still obtains a smaller tail gap ($6.47\times10^{-3}$), although its transient cost remains substantially larger than EF-TopK.

Thus Experiment~11B passes the communication-aware feasibility gate: there exists a nontrivial calibration regime that improves on the selected conventional sparse baseline even after dense calibration traffic is charged.

\subsection{Calibration traffic is now the dominant bottleneck}

The positive result comes with a severe systems limitation. At $C=25$,
\begin{equation}
  \frac{B_{\rm cal}}{B_{\rm total}}
  =
  \frac{518400}{533597}
  \approx97\%.
\end{equation}
At $C=75$, calibration still accounts for about 92\% of payload. The sparse event stream remains cheap; recurring trustworthy magnitude information is now the dominant communication cost.

The bottleneck has therefore moved from ``how to size a spike'' to ``how to maintain a sufficiently tight global descent certificate without repeatedly transmitting dense gradient vectors.''

\subsection{Naive stale gradients are unsafe}

Ignoring the drift radius and treating the stale calibration center as the current gradient is catastrophically unstable. In the controlled sweeps, naive rules accumulate very large objective errors and a majority of their accepted events can be globally harmful. This control demonstrates that occasional dense gradients alone are not the mechanism: the uncertainty envelope is essential.

\subsection{Certificate-triggered refresh does not yet save meaningful bits}

The adaptive-refresh variant recalibrates only after an uncertified event and subject to a minimum refresh interval. In the present native rotated benchmark it converges toward nearly the same refresh frequency as the corresponding periodic schedule once the uncertainty set has expanded. It therefore produces little communication reduction. This suggests that a useful adaptive policy must decide \emph{which part} of the certificate needs refresh, not merely replace a fixed clock by a binary refresh request.

\subsection{First-passage tightening remains insufficient}

The firing-time estimator from Experiments~10C and 11A was used to refresh only the firing client's contribution between dense calibrations. With a 99\% empirical signed-gradient error margin, the rule admits nonzero harmful-event rates and substantially distorts slow-client participation. Tightening the margin to 99.9\% reduces some unsafe activity but also makes the rule much more conservative. It still does not outperform the deterministic periodic certificate.

Representative 99.9\% results include a diagonal tail gap around $6.5\times10^{-2}$ where the deterministic certificate reaches numerical zero. In the rotated $C=100$ case the timing-assisted rule remains above $0.1$ tail gap. First-passage information should therefore remain a diagnostic signal, not part of the core certificate at this stage.

\subsection{Representation geometry and certificate geometry are coupled}

The oracle basis-aligned control provides the strongest mechanistic result. In the known common eigenbasis from Experiment~10D, both event coding and gradient-drift certification become coordinate aligned. With only three calibrations ($C=1000$), the rotated problem returns to
\begin{equation}
  J_{\rm tail}\approx0,
  \qquad
  J_{\rm whole}=0.209,
  \qquad
  B=27056\;\text{bits}.
\end{equation}
This is essentially the diagonal calibrated-certificate result at the same communication cost.

Therefore the high calibration demand of native rotated Experiment~11B is not intrinsic to periodic calibration. It is caused by the interaction between sparse coordinate events and a rapidly expanding uncertainty set in a misaligned representation:
\begin{equation}
  \boxed{\text{event representation and certificate representation should be designed together.}}
\end{equation}
The eigenbasis remains an oracle rather than a deployable neural-network solution, but it strongly motivates later block, preconditioned, or low-rank communication coordinates.

\subsection{Safety and client participation}

The deterministic periodic certificate does not obtain safety simply by silencing the slowest clients. In the rotated transition sweep, the two $T=20$ clients generate approximately 19\% of candidate events. Their accepted-event share is approximately 17--19\% at the useful $C=25$--$75$ operating points, reasonably close to their intended combined weight of 20\%. This is not perfect fairness, but it is far less distorted than the hard competition studied in Experiment~10B.

\subsection{Experiment 11B verdict}

\begin{equation}
  \boxed{\text{Periodic calibrated descent certificate: PASS}}
\end{equation}
There exist calibration periods that improve optimization substantially while retaining a communication advantage after dense calibration bits are counted.

\begin{equation}
  \boxed{\text{Stale calibration without uncertainty bound: FAIL}}
\end{equation}
The naive rule is unsafe and can become unstable.

\begin{equation}
  \boxed{\text{Sparse native certification under rotated geometry: PARTIAL PASS}}
\end{equation}
Useful Pareto points exist, but tight guarantees require much more frequent calibration than in the aligned problem.

\begin{equation}
  \boxed{\text{First-passage timing as certificate tightening: FAIL for now}}
\end{equation}
Empirical timing margins either sacrifice safety or become too conservative.

\begin{equation}
  \boxed{\text{Basis-aligned calibrated certificate: strong ORACLE PASS}}
\end{equation}
A good shared representation simultaneously restores sparse event efficiency and sparse certificate maintenance.

The principal conclusion is
\begin{equation}
  \boxed{\textbf{occasional trustworthy magnitude information can make descent-aware sparse events practical, but calibration payload and representation geometry are now the dominant design problems.}}
\end{equation}
